Je remercie sincèrement mon encadrant académique, Dr. Riheb Naili, pour son suivi et pour les remarques
qu'il m'a adressées tout au long de ce projet de fin d'études. Ses conseils m'ont souvent poussé à reprendre
certains choix, à mieux justifier la démarche suivie et à présenter PharmacieConnect d'une manière plus claire.

Je remercie également \textbf{IT-STORM Consulting}, mon organisme d'accueil, ainsi que mon encadrant
professionnel, \textbf{Achraf BCHIRI}, pour l'accueil et les orientations donnés pendant le stage externe.
Dans ce cadre, j'ai pu relier le développement réalisé à un besoin concret : permettre à un citoyen de
trouver plus facilement une pharmacie de garde, surtout lorsque l'information doit être consultée rapidement.

Je remercie l'équipe pédagogique de l'Institut Supérieur d'Informatique et de Multimédia de Gabès pour
l'accompagnement assuré durant cette période. Les échanges avec les enseignants m'ont aidé, dans notre cas,
à garder un lien entre les choix techniques, les objectifs du mémoire et les exigences de la formation.

Je souhaite aussi exprimer ma reconnaissance aux enseignants qui ont contribué à ma formation. Les notions
acquises en développement web, en développement mobile, en bases de données, en sécurité applicative et en
génie logiciel ont été utiles à chaque étape du projet, parfois de manière très directe, parfois comme repères
pour prendre de meilleures décisions.

Enfin, je remercie ma famille et mes amis pour leur soutien. Leur présence m'a beaucoup aidé à continuer le
travail avec régularité, y compris pendant les moments plus difficiles de l'analyse, du développement, des
tests et de la rédaction.
